\section{IMPLEMENTATION}

\toolname is implemented as a Selenium-based fuzzing pipeline that
generates HTML/CSS test cases, runs them in a target browser, detects
under-invalidation bugs, and post-processes bug reports.

\subsection{Fuzzing Process Execution}
\begin{itemize}
    \item \toolname generates a random DOM tree with base CSS
    styles and modified CSS styles.
    \item The generated page is executed in a real browser through
    Selenium WebDriver.
    \item The harness compares layout after incremental style changes
    against layout after a fresh reload to detect under-invalidation
    bugs.
\end{itemize}

\subsection{Test Case Mutation and Trimming}
\begin{itemize}
    \item When a bug is found, \toolname repeatedly applies
    minimization steps and keeps only changes that still reproduce the
    bug.
    \item The minimizer removes unnecessary elements, CSS rules, and
    style changes from the generated test case.
    \item \texttt{Minify\_MoveStyleChangeToFirstLoad} moves nonessential
    modified styles into the base styles, leaving the modified styles
    most critical to reproducing the bug.
    \item The final minimized page is saved with metadata describing the
    affected styles and layout differences.
\end{itemize}

\subsection{Bug Categorization and Deduplication}
\begin{itemize}
    \item \toolname converts minimized bugs into structural
    representations that include the DOM tree, base styles, and modified
    styles.
    \item The system compares new bugs against existing bug groups to
    identify repeated instances of the same underlying issue.
    \item Matching bugs are placed into existing groups, while new
    patterns can form new bug groups with generated rules.
\end{itemize}

\subsection{Supported Software and Browser Platforms}
\begin{itemize}
    \item \toolname supports browser testing through Selenium
    variants for Chrome, Firefox, and Safari.
    \item Browser variants can be configured with custom binaries,
    WebDriver paths, viewport sizes, and headless execution.
    \item Bug reports are grouped using structural rules so repeated
    instances of the same underlying bug can be categorized together.
\end{itemize}
